\documentclass[11pt]{article}
\newcommand{\shorttitle}{Simultaneous torsion specialization}
\usepackage[T1]{fontenc}
\usepackage{lmodern,amsmath,amssymb,amsthm,mathtools,booktabs,microtype}
\usepackage[letterpaper,left=1.03in,right=1.03in,top=.90in,bottom=.90in]{geometry}
\usepackage{xcolor,fancyhdr,enumitem}
\usepackage{hyperref}
\definecolor{accent}{HTML}{1D5660}
\definecolor{ink}{HTML}{192E38}
\hypersetup{colorlinks=true,linkcolor=accent,urlcolor=accent,citecolor=accent}
\newtheorem{theorem}{Theorem}[section]
\newtheorem{corollary}[theorem]{Corollary}
\newtheorem{proposition}[theorem]{Proposition}
\setcounter{tocdepth}{1}
\setlength{\headheight}{14pt}
\setlength{\emergencystretch}{1.3em}
\setlength{\parindent}{1.15em}
\setlist{itemsep=0.15em,topsep=0.4em}
\pagestyle{fancy}\fancyhf{}
\fancyhead[L]{\small\color{accent}\shorttitle}
\fancyhead[R]{\small\color{accent}25 November 2026}
\fancyfoot[C]{\small\thepage}
\renewcommand{\headrulewidth}{0.3pt}
\newcommand{\paperheading}[3]{\noindent{\sffamily\small\bfseries\color{accent}#1}\par\vspace{.75em}
  {\raggedright\sffamily\bfseries\LARGE\color{ink}#2\par}\vspace{.6em}
  {\noindent\large #3\par}\vspace{1em}}

\hypersetup{pdftitle={The exact cost of obstructing simultaneous torsion specialization}}
\begin{document}
\paperheading{ARITHMETIC OF ROOTS OF UNITY}{How far must we search for a simultaneous $p$-adic twist?}{An exact answer to the packet-size question in the ProveIt quantitative-twists report}
\noindent
A step in the Lean proof of Dirichlet GRH chooses roots of unity where several integral polynomials keep small valuations. A previous analysis asked whether one can bound the required order exponent by a polynomial in the total valuation budget and inverse tolerance, uniformly over the number of polynomials, when one bounds the \emph{sum of the positive valuation increases}. We determine the exact worst-case answer in one variable. With an unrestricted coefficient field and fixed positive tolerance $\delta$, the least guaranteed exponent has logarithm $B/\delta+O(1)$ as the budget $B$ grows. Balanced sets of roots of unity give examples that attain the optimum. For a fixed coefficient field of absolute ramification index $e$, the same problem admits a linear bound and even permits zero loss. For a cyclotomic coefficient field we also determine the exact answer at every tolerance. The proofs are independent of RH.
\section{The step in the Lean proof}
The five public theorems at commit \texttt{72fead102f66} concern the Riemann zeta function and Dirichlet $L$-functions. They include zero-freeness for real part greater than $1/2$, subject to the exclusion for the principal character at $s=1$, and the critical-line statement throughout the open strip. The readable and compact editions have independent machine-check records for those statements. We used the pinned sources and Andrew Sutherland's verification record. We did not run a new Lean build.\footnote{\href{https://github.com/lsernau/grh-lean/tree/72fead102f663d5c7a407964a13dc0cf9d977207}{Upstream sources} and \href{https://math.mit.edu/~drew/shortgrh/}{independent verification record}. Exact source paths and prior-work references are in \texttt{sources.md}.}
\par
In the readable edition, short integral forms and an index bound, a supply of independent translates, and a lower bound on the corrected covolume of integral forms enter \texttt{star\_inequality}. Comparing the upper and lower bounds leaves a one-sided estimate on a smoothed prime sum. The construction on the lower-bound side chooses anticyclotomic characters so that several Selmer and special-value quantities are small at once. The local input in \texttt{RootsOfUnity.lean} selects specializations of polynomials by bounding how much valuation can be concentrated on a given shell of torsion points. Its order bound is chosen before the polynomials are supplied. This quantifier is the origin of the extremal problem here.
\par
The earlier quantitative-twists report gave a polynomial upper bound for a \emph{fixed} finite number of constraints. Another incoming report gave a bound independent of their number by exploiting a divergent harmonic sum. The question concerns the price of uniformity. Section~\ref{sec:torsion-minimax} answers it for every finite search range. Section~\ref{sec:limit} gives the transition as the tolerance approaches zero. Section~\ref{sec:field} applies the result to a fixed coefficient field and extends zero-loss selection to several variables.
\section{The exact cost of obstructing simultaneous torsion specialization}
\label{sec:torsion-minimax}

We answer Question (ii) of the quantitative-twists report by working in one variable. A uniform polynomial bound in all dimensions would have to hold here. Logarithms are natural unless a base is specified.

Fix a prime $p$, a radial lower index $h_0\ge1$ and $\rho\ge
1/\varphi(p^{h_0})$. Put $N_j=(p-1)p^{j-1}$, $r_j=1/N_j$, and use
$v(p)=1$, $v(0)=+\infty$. The $h$-th shell consists of the points
$z=\zeta-1$ for $\zeta$ of exact order $p^h$. Consider a finite packet of
nonzero polynomials $g_i\in\mathcal O_{\mathbf C_p}[T]$ with witnesses
$x_i$ satisfying $v(x_i)\ge\rho$ and $t_i=v(g_i(x_i))<\infty$. Its
budget is $B=\sum_i t_i$. Its loss at $z$ is
\[
             D(z)=\sum_i\max\{0,v(g_i(z))-t_i\}.
\]
A shell is blocked at tolerance $\delta\ge0$ if every point on it has
$D(z)>\delta$. Define the extended-real function
\[
 S_h(L)=\frac1{N_h}\sum_{j=1}^{h-1}N_j\left\lfloor\frac{L}{N_j}\right\rfloor
 \quad(0\le L<N_h),\qquad S_h(N_h)=+\infty .
\]
Let $m_h$ be the least integer $L\in\{0,\ldots,N_h\}$ for which $S_h(L)>\delta$.
In particular the empty sum at $h=1$ is zero. The strict threshold is
essential when $\delta$ is an attained loss.

\begin{theorem}[Exact obstruction budget]\label{thm:exact-obstruction}
The minimum budget of a finite packet blocking every shell from $h_0$
through $H$ is
\[
              C_{p,\delta}(h_0,H)=\sum_{h=h_0}^{H}\frac{m_h}{N_h}.
\]
This minimum is attained using only monic linear polynomials with a common
witness $0$ and coefficients in the ring of integers of
$\mathbf Q_p(\zeta_{p^H})$. It also holds if $t_i$ is replaced throughout by
the Gauss valuation $G_\rho(g_i)=\min_k(v([T^k]g_i)+k\rho)$.
\end{theorem}

\begin{proof}
For $g=A\prod_\alpha(T-\alpha)$, count roots with multiplicity. Multiplicativity
of the Gauss valuation gives $G_s(g)=v(A)+\sum_\alpha\min(s,v(\alpha))$.
This follows by multiplying the reductions of scaled polynomials at rational
weights, or by choosing the smallest exponent attaining each coefficient
minimum. Integral coefficients give $G_0(g)\ge0$. Therefore
\[
 \sum_{h=h_0}^H r_h\#\{\alpha:v(\alpha)=r_h\}
 \le G_\rho(g)-G_0(g)\le G_\rho(g)\le v(g(x_i))
\]
when $g=g_i$. On the $h$-th shell every root with valuation different
from $r_h$ contributes exactly its Gauss term, by the ultrametric inequality.
Thus for every $z$ on this shell,
\[
 D(z)\le\sum_{\substack{\alpha\text{ across the packet}\\v(\alpha)=r_h}}
                       (v(z-\alpha)-r_h).
\]
Indeed $G_{r_h}(g_i)\le G_\rho(g_i)\le t_i$ and every summand on the
right is nonnegative. This argument works with either stated baseline.

Fix a primitive $p^h$-th root $\xi$. For each of the $L$ roots $\alpha$ on
this shell mark an exponent $b$ such that $\xi^b-1$ is closest to $\alpha$.
Here $b$ is a unit modulo $p^h$, and ties may be resolved arbitrarily. If
$a\ne b$ modulo $p^h$, the nearest-point choice implies
$v(\xi^a-1-\alpha)\le v(\xi^a-\xi^b)$. Writing $k=v_p(a-b)$ gives
\[
 v(\xi^a-\xi^b)-r_h
 =\frac{p^k-1}{N_h}=\frac1{N_h}\sum_{j=1}^k N_j.
\]
The first equality uses $v(\xi^{a-b}-1)=1/N_{h-k}$, from the cyclotomic
polynomial at $1$. Marks are allowed to coincide.

If $L<N_h$, begin with a least occupied unit residue class modulo $p$ and
at every stage choose its least occupied lift modulo the next power of $p$.
There are $p-1$ choices at the first stage and $p$ at every subsequent one.
The resulting class modulo $p^j$ contains at most $\lfloor L/N_j\rfloor$
marks, by induction. At level $h$ it is empty. Summing the displayed distance
identity over marks at this unmarked endpoint $a$ gives
$D(\xi^a-1)\le S_h(L)$. Consequently every blocked shell has at least
$m_h$ roots. The Gauss-budget inequality, summed over all polynomials,
establishes the lower bound in the theorem.

For equality, on level $h$ mark the first $m_h$ positive integers prime to
$p$ in their natural order. For each marked $b$, insert the linear
polynomial $T-(\xi^b-1)$ in the packet. Its valuation at $0$ and its
$\rho$-Gauss valuation are both $r_h$. In successive blocks of $N_j$ unit
exponents, every unit residue class modulo $p^j$ appears exactly once.
Hence every residue class at level $j$ contains at least
$\lfloor m_h/N_j\rfloor$ marks. A marked endpoint gives an infinite loss.
At any unmarked endpoint the distance identity gives loss at least
$S_h(m_h)>\delta$. This also proves the assertion when $m_h=N_h$ since all
endpoints are then marked. Linear factors from other shells have zero
positive excess there. Taking all these polynomials gives exactly the
claimed cost.
\end{proof}

At each shell one may multiply its linear factors into one polynomial,
since on that shell each individual excess is nonnegative and on every
other shell their excesses have the same sign. Thus at most
$H-h_0+1$ polynomials suffice if degree is unrestricted.
The balanced residue counts are an instance of the fixed-divisor phenomenon
studied in Bhargava's $P$-orderings (M.~Bhargava, \emph{J. Reine Angew.
Math.} 490 (1997), 101--127, \href{https://doi.org/10.1515/crll.1997.490.101}{doi:10.1515/crll.1997.490.101}).
Here the positive-part loss and distribution of the witness budget across
different cyclotomic shells determine the global optimum.

\begin{corollary}[Exponential depth is necessary]\label{cor:exponential-depth}
For fixed $p,h_0$ and $\delta>0$,
\[
 C_{p,\delta}(h_0,H)=\delta\log H+O_{p,h_0,\delta}(1).
\]
The least universal upper shell index given budget $B$ is the first $H$
with $C_{p,\delta}(h_0,H)>B$. Consequently its logarithm equals
$B/\delta+O_{p,h_0,\delta}(1)$ as $B\to\infty$.
\end{corollary}
\begin{proof}
The upper estimate $S_h(L)\le(h-1)L/N_h$ holds for $L<N_h$. For large
$h$ it follows that $m_h/N_h\ge\delta/(h-1)$, also when $m_h=N_h$.
For the reverse bound let $q=h-\lceil 3\log_p h\rceil$; eventually
$1\le q<h$. For any $L<N_h$,
\[
 S_h(L)\ge\frac{qL-\sum_{j=1}^q N_j}{N_h}
          =\frac{qL-p^q+1}{N_h}.
\]
Take $L=\lfloor(\delta N_h+p^q)/q\rfloor+1$, eventually smaller than
$N_h$. This makes $S_h(L)>\delta$. Since $p^q/N_h=O_p(h^{-3})$ and
$h-q=O_p(\log h)$, one gets
$m_h/N_h\le\delta/h+O_{p,\delta}(\log(h)/h^2)$. Summation and convergence
of the error series give the formula; the strict inequality in the exact
criterion gives the asserted least index.
\end{proof}

\paragraph{Zero tolerance.}
For $h\ge2$, $S_h(L)$ first becomes positive at $L=p-1$; the same number
marks the whole first shell. Thus the exact finite obstruction cost at
$\delta=0$ is $\sum_{h=h_0}^{H}p^{1-h}$. For example, at $p=2$ and $h_0=1$,
arbitrarily many shells can be blocked at zero tolerance using budget less
than $2$. This shows why the positive-tolerance hypothesis in the asymptotic
law matters.


\section{The transition at small tolerance}\label{sec:limit}
Write $a_h(\delta)=m_h/N_h$ and let $H_*(B,\delta)$ be the first $H\ge h_0$ with $C_{p,\delta}(h_0,H)>B$. This index bounds the \emph{exponent}: the order of the root is at most $p^{H_*}$. The comparison is strict because the extremum in Theorem~\ref{thm:exact-obstruction} is attained. The formula is readily computable by binary search for $m_h$ with integer or rational arithmetic.

\paragraph{A finite example.}
At $p=2$, $h_0=1$ and $\delta=1/4$ the optimal costs are as follows.
\begin{center}\begin{tabular}{@{}lrrrrrr@{}}\toprule
Shell exponent $h$ & 1 & 2 & 3 & 4 & 5 & 6\\
\midrule
Primitive roots $N_h$ & 1&2&4&8&16&32\\
Minimum marks $m_h$ &1&1&2&2&3&4\\
Cumulative budget & $1$ & $3/2$ & $2$ & $9/4$ & $39/16$ & $41/16$\\
\bottomrule\end{tabular}\end{center}
A budget of exactly $2$ can block orders $2,4,8$. Every packet of that budget admits a successful point of order at most $16$. At order $8$, one mark attains minimum loss exactly $1/4$ across the primitive roots; the definition counts equality as success, so blocking this shell requires a second mark.

\begin{theorem}[The fixed-budget transition]\label{thm:transition}
Set $s_h=p^{1-h}$ and
\[
 B_0=\sum_{h=h_0}^\infty s_h=\frac{p^{2-h_0}}{p-1}.
\]
For fixed $p,h_0,B\ge0$ (with $p$ prime and $h_0\ge1$), the least guaranteed exponent as $\delta\downarrow0$ obeys the following three laws.
\begin{enumerate}
\item If $B<B_0$, it eventually equals the first $H$ with $\sum_{h=h_0}^H s_h>B$.
\item If $B=B_0$, put $t=\lceil1+\log_p(1/\delta)\rceil$. For all sufficiently small positive $\delta$, the answer is $t$ for odd $p$ and $t+1$ for $p=2$.
\item If $B>B_0$, then $\displaystyle\lim_{\delta\downarrow0}\delta\log H_*(B,\delta)=B-B_0$.
\end{enumerate}
\end{theorem}
\begin{proof}
For each fixed shell, $a_h(\delta)=s_h$ for every sufficiently small $\delta\ge0$. Indeed for $h\ge2$ the first nonzero finite score occurs at $L=p-1$, where $S_h(p-1)=s_h$. For $h=1$ all leaves must be marked. This proves the subcritical case and shows that $H_*(B,\delta)\to\infty$ when $B\ge B_0$.

The estimates in Corollary~\ref{cor:exponential-depth} give the uniform summable bound, for $h$ sufficiently large and $0\le\delta\le1$,
\[
 0\le a_h(\delta)-\delta/h
 \ll_p \frac{\log(2h)}{h^2}.
\]
For clarity, the upper bound uses $q=h-\lceil3\log_p h\rceil$ and
$a_h\le \delta/q +p^q/(qN_h)+1/N_h$. The lower bound follows from $S_h(L)\le(h-1)L/N_h$ for $L<N_h$, and also holds if $m_h=N_h$. Dominated convergence on the positive integers now yields, along \emph{any} $H\to\infty$, $\delta\downarrow0$,
\[
 C_{p,\delta}(h_0,H)-\delta\log H\longrightarrow B_0.
\]
Here $\sum_{h=h_0}^H 1/h-\log H$ is bounded. At $H=H_*$ the overshoot above $B$ is at most $a_H(\delta)$ once $H>h_0$. The uniform bound shows that this tends to zero, proving the supercritical case.

For the critical case take $\delta$ small enough that $t>\max(2,h_0)$. All earlier costs equal $s_h$. At $h=t$ the tolerance is in $[s_t,p s_t)$. All $N_j$ are multiples of $p-1$, so the score is constant between consecutive multiples of $p-1$ in the mark count. Hence $a_t\ge2s_t$. The remaining baseline tail is $\sum_{h>t}s_h=s_t/(p-1)$. For $p>2$ the extra cost already exceeds this tail and $H_*=t$. For $p=2$, the exact scores at $L=1,2$ are $s_t,4s_t$ since $t\ge3$. Thus $a_t=2s_t$, the cumulative cost is exactly $B_0$, and one further positive-cost shell makes it strictly larger.
\end{proof}

For $p=2$ and $h_0=1$, the critical budget is $2$. Raising the budget from exactly $2$ to $3$ changes the asymptotic exponent from $\lceil1+\log_2(1/\delta)\rceil+1$ to $\exp((1+o(1))/\delta)$. Thus the failure of a packet-size-independent polynomial bound already occurs with \emph{fixed} budget and arbitrarily small tolerance.

\section{Zero loss over a fixed coefficient field}\label{sec:field}
\noindent
Suppose all the polynomials are defined over a finite extension
$K/\mathbf Q_p$ of absolute ramification index $e$. For an irreducible
factor of degree $d$ containing a root of valuation $1/N_h$, all its roots
have that valuation, by invariance of the unique extended valuation.
The norm has valuation $d/N_h\in e^{-1}\mathbf Z$. A nonempty shell thus
costs at least $1/\gcd(e,N_h)\ge1/e$. Among any $\lfloor eB\rfloor+1$
eligible shells there is one with no root at its radius. All primitive
evaluations there are simultaneously bounded above by their witness
valuations. 


\subsection*{An exact formula over a cyclotomic coefficient field}
Fix $s\ge1$, a compatible system of primitive roots $\xi_h$ with $\xi_{h+1}^p=\xi_h$, and the coefficient field $K_s=\mathbf Q_p(\xi_s)$. Its absolute ramification index is $N_s$. On an arbitrarily high shell, a packet over this field has precisely the same optimum as a packet on shell $s$.

\begin{theorem}[Descent to the coefficient level]\label{thm:descent}
For packets with coefficients in $\mathcal O_{K_s}$, the exact minimum witness budget needed to block the entire interval $[h_0,H]$ at tolerance $\delta\ge0$ is
\[
       \sum_{h=h_0}^H a_{\min(h,s)}(\delta),
       \qquad a_j(\delta)=m_j/N_j .
\]
For $h_0\ge s$ the least guaranteed order exponent is exactly
\[
        h_0+\left\lfloor B/a_s(\delta)\right\rfloor .
\]
\end{theorem}
\begin{proof}
For $h\le s$ both the lower bound and the optimizing linear factors of Theorem~\ref{thm:exact-obstruction} already apply. Suppose $h>s$. Partition all roots of valuation $r_h$ in the packet into conjugacy orbits over $K_s$, repeated according to multiplicity. If an orbit has degree $d$, its norm valuation $d/N_h$ lies in $N_s^{-1}\mathbf Z$, so $u=dN_s/N_h$ is a positive integer.

Pick a member $\alpha$ and a nearest point $\xi_h^b-1$. Every automorphism fixing $K_s$ acts on $\xi_h$ by an exponent congruent to $1$ modulo $p^s$. Consequently each root in that orbit has a nearest point whose exponent is congruent to $b$ modulo $p^s$. Ties cause no problem: one conjugating automorphism can be chosen for each member. Assign $u$ marks to the single unit class $b\bmod p^s$. If a test class $a\bmod p^s$ is unmarked, its $p$-adic distance in exponent from each of the nearest exponents for this orbit is given by the same $k=v_p(a-b)<s$. At every lift of $a$ to shell $h$, this orbit contributes to the upper bound for the loss at most
\[
             d\,\frac{p^k-1}{N_h}=u\,\frac{p^k-1}{N_s}.
\]
Write $U=\sum u$. The actual root budget on shell $h$ is $\sum d/N_h=U/N_s$. If $U<m_s\le N_s$, the least-occupied path on shell $s$ supplies an unmarked class with total loss at most $S_s(U)\le\delta$ at every lift. Thus blocking requires a cost at least $a_s$ on each shell $h>s$.

Conversely choose the first $m_s$ positive unit exponents $b$ at level $s$. At level $h>s$ take one polynomial for each such $b$,
\[
            F_{h,b}(T)=(1+T)^{p^{h-s}}-\xi_s^b.
\]
Its roots are exactly the $\xi_h^c-1$ with $c\equiv b\pmod {p^s}$. Its value at the origin and its Gauss baseline both have valuation $r_s$. At $T=\xi_h^a-1$, its value is $\xi_s^a-\xi_s^b$. Hence the entire loss profile at level $s$ pulls back to level $h$. Values on all other shells have zero positive excess by the root factorization. Combining these polynomials with the linear constructions below level $s$ realizes the formula with a common witness $0$. The last assertion follows from strict budget crossing.
\end{proof}
For example, take $p=2$, $s=3$ and $\delta=1/4$. At every $h\ge3$ the exact price of blocking the shell over $\mathbf Q_2(\zeta_8)$ is $a_3=1/2$. A field that already contains all $2^h$-th roots permits the smaller costs in the preceding table as $h$ grows. This identifies exactly where extra coefficient ramification can aid an adversarial packet.


\subsection*{Several variables}
The fixed-field result does not require witnesses in the coefficient field. Let $K/\mathbf Q_p$ be finite of absolute ramification index $e$ and take a finite family $F_i\in\mathcal O_K[T_1,\dots,T_s]\setminus\{0\}$. Write
\[
 A_i=\min_{\mathbf a}\{v([T^{\mathbf a}]F_i)+\rho|\mathbf a|\},
 \qquad \sum_i A_i\le B .
\]
Then there is a single tuple $z_j=\zeta_j-1$, each of order exponent $h_j\ge h_0$ for $\zeta_j$, such that
\[
                 v(F_i(z_1,\ldots,z_s))\le A_i\quad\text{for every }i.
\]
This gives zero positive excess relative to all witnesses of coordinate valuations at least $\rho$. Explicit exponents $h_j\le M_j$ are available before the polynomials are given:
\[
 M_1=h_0+\lfloor eB\rfloor,\qquad
 M_{j+1}=h_0+\lfloor e N_{M_j}B\rfloor.
\]

\begin{proof}
At a given step expand a partially specialized polynomial as
$\sum_{\nu} g_{i,\nu}(T)Y^\nu$, where $Y$ denotes all later variables. Its multivariate Gauss valuation $A_i'$ is $\min_\nu(G_\rho(g_{i,\nu})+\rho|\nu|)$. Choose one nonzero coefficient polynomial $g_i$ attaining that minimum. Then $\sum_i G_\rho(g_i)\le\sum_i A_i'\le B$, initially by hypothesis. Apply the one-variable fixed-field result with these Gauss baselines, which the root-budget proof also permits. At a common point $z$ one has $v(g_i(z))\le G_\rho(g_i)$ for every $i$. The new Gauss valuation after specialization is at most $v(g_i(z))+\rho|\nu|\le A_i'$. This also ensures every new polynomial is nonzero and integral. Iterate over all coordinates.

The coefficient field after $j$ steps is contained in $K(\mu_{p^{M_j}})$, as the stated bounds are nondecreasing and the cyclotomic levels are nested. Its absolute ramification index is at most $e[K(\mu_{p^{M_j}}):K]\le eN_{M_j}$. The one-variable count of root-bearing shells gives the recurrence. At the final coordinate the Gauss valuation of the resulting constant is its valuation.
\end{proof}

\paragraph{What this gives for arithmetic twist families.}
This bound can be used whenever the coefficient field is controlled \emph{before} choosing the order cutoff. The worst-case packets in Theorem~\ref{thm:exact-obstruction} live in finite cyclotomic fields whose ramification grows with the largest forbidden shell. An application to a particular Selmer or special-value family therefore turns on its coefficient fields, or on structural restrictions on its interpolation polynomials. The extremal packets are explicit integral polynomials; their occurrence as Selmer determinants would be an additional arithmetic theorem.

\paragraph{Further research questions.}
Determine the optimum at fixed packet cardinality between the single-polynomial and unrestricted regimes. Determine the exact obstruction cost over a finite extension that is not cyclotomic. In several variables, is the iterated growth of the displayed fixed-field bound necessary, or can one use the common Galois action on the coordinates to reduce the search? The independent analytic applications in \texttt{prime-sampling.pdf} pursue zero location and simplicity by a different part of the same Lean proof.

\end{document}
